% GENERATED FILE. Edit canonical lessons, then run npm run generate:books.
% canonical-source-hash: fnv1a64:055c476a66627f28
% canonical-lessons: ML-C448-pratisandhi, ML-C448-velluvili, ML-C448-tatassam, ML-C448-visakalanam-ceyyuka, ML-C448-riti

\chapter{Facing a Crisis}
\label{ch:ml-facing-a-crisis}
\hlchaptermodality{448}

\begin{chapteropening}
\textbf{By the end of this chapter:} \emph{I can say a crisis, a challenge, an obstacle, to analyse and a method.}

Complete the last lesson of chapter 448: I can say a crisis, a challenge, an obstacle, to analyse and a method.
\end{chapteropening}

\section[\texorpdfstring{\ml{പ്രതിസന്ധി} (pratisandhi)}{പ്രതിസന്ധി (pratisandhi)}]{\ml{പ്രതിസന്ധി} (pratisandhi) — a crisis}

\label{lesson:ML-C448-pratisandhi}



\begin{quote}
Before the new one: say the Malayalam for to aim at, then the Malayalam for efficiency.
\end{quote}

\subsection*{You'll want to know: \ml{പ്രതിസന്ധി}}
\textbf{\ml{പ്രതിസന്ധി}} — \emph{pratisandhi} — \textquotedblleft{}a crisis\textquotedblright{}.

\begin{quote}
\ml{കമ്പനി} \ml{ഇപ്പോൾ} \ml{വലിയ} \ml{പ്രതിസന്ധിയിലാണ്}.
\end{quote}

\begin{quote}
\emph{kampani ippōḷ valiya pratisandhiyilāṇŭ.}
\end{quote}

\begin{quote}
The company is now in a big crisis.
\end{quote}

In a crisis: \ml{പ്രതിസന്ധിയിൽ}.

\begin{cousinweb}
\textbf{\ml{പ്രതിസന്ധി}} (\emph{pratisandhi}) is built from Sanskrit \emph{prati-}, \textquotedblleft{}against\textquotedblright{}, and \emph{saṁdhi}, \textquotedblleft{}a junction, a joint\textquotedblright{} (Monier-Williams). Monier-Williams gives \emph{pratisaṁdhi} as \textquotedblleft{}reunion\textquotedblright{}; the sense \textquotedblleft{}a crisis\textquotedblright{} is the modern one.
\end{cousinweb}

\begin{grammarlens}[title={What you've built}]
One more word for saying what happens if.
\end{grammarlens}

\subsection*{Your turn}
\begin{itemize}
\raggedright
  \item \emph{Say it:} \emph{pratisandhi}
  \item \emph{Say it:} \emph{pratisandhi}, once more
  \item \emph{Say it:} \emph{kampani ippōḷ valiya pratisandhiyilāṇŭ.} — The company is now in a big crisis.
  \item \emph{Recall:} say the Malayalam for to aim at, then the Malayalam for efficiency, then say \emph{pratisandhi} again
\end{itemize}

\begin{tcolorbox}[breakable,colback=teal!4,colframe=teal!35!black,title={Before you move on}]
What does \ml{പ്രതിസന്ധി} mean? (\textquotedblleft{}A crisis\textquotedblright{}.) Say it once more.
\end{tcolorbox}
\section[\texorpdfstring{\ml{വെല്ലുവിളി} (velluviḷi)}{വെല്ലുവിളി (velluviḷi)}]{\ml{വെല്ലുവിളി} (velluviḷi) — a challenge}

\label{lesson:ML-C448-velluvili}



\begin{quote}
Before the new one: say the Malayalam for efficiency, then the Malayalam for a crisis.
\end{quote}

\subsection*{You'll want to know: \ml{വെല്ലുവിളി}}
\textbf{\ml{വെല്ലുവിളി}} — \emph{velluviḷi} — \textquotedblleft{}a challenge\textquotedblright{}.

\begin{quote}
\ml{പുതിയ} \ml{ജോലി} \ml{ഒരു} \ml{വലിയ} \ml{വെല്ലുവിളിയാണ്}.
\end{quote}

\begin{quote}
\emph{putiya jōli oru valiya velluviḷiyāṇŭ.}
\end{quote}

\begin{quote}
The new job is a big challenge.
\end{quote}

To challenge is \ml{വെല്ലുവിളിക്കുക}.

\begin{cousinweb}
\textbf{\ml{വെല്ലുവിളി}} (\emph{velluviḷi}) is a native Malayalam compound: \ml{വെല്ലുക}, \textquotedblleft{}to overcome, to surpass\textquotedblright{} (DEDR 5493), and \ml{വിളി}, \textquotedblleft{}a call\textquotedblright{} (DEDR 5433): a call to see who wins.
\end{cousinweb}

\begin{grammarlens}[title={What you've built}]
One more word for saying what happens if.
\end{grammarlens}

\subsection*{Your turn}
\begin{itemize}
\raggedright
  \item \emph{Say it:} \emph{velluviḷi}
  \item \emph{Say it:} \emph{velluviḷi}, once more
  \item \emph{Say it:} \emph{putiya jōli oru valiya velluviḷiyāṇŭ.} — The new job is a big challenge.
  \item \emph{Recall:} say the Malayalam for efficiency, then the Malayalam for a crisis, then say \emph{velluviḷi} again
\end{itemize}

\begin{tcolorbox}[breakable,colback=teal!4,colframe=teal!35!black,title={Before you move on}]
What does \ml{വെല്ലുവിളി} mean? (\textquotedblleft{}A challenge\textquotedblright{}.) Say it once more.
\end{tcolorbox}
\section[\texorpdfstring{\ml{തടസ്സം} (taṭassaṁ)}{തടസ്സം (taṭassaṁ)}]{\ml{തടസ്സം} (taṭassaṁ) — an obstacle}

\label{lesson:ML-C448-tatassam}



\begin{quote}
Before the new one: say the Malayalam for a crisis, then the Malayalam for a challenge.
\end{quote}

\subsection*{You'll want to know: \ml{തടസ്സം}}
\textbf{\ml{തടസ്സം}} — \emph{taṭassaṁ} — \textquotedblleft{}an obstacle\textquotedblright{}.

\begin{quote}
\ml{മഴ} \ml{കളിക്ക്} \ml{തടസ്സമായി}.
\end{quote}

\begin{quote}
\emph{maḻa kaḷikkŭ taṭassamāyi.}
\end{quote}

\begin{quote}
The rain became an obstacle to the game.
\end{quote}

To become an obstacle to: -\ml{ക്ക്} \ml{തടസ്സമാകുക}.

\begin{cousinweb}
\textbf{\ml{തടസ്സം}} (\emph{taṭassaṁ}) is usually traced to native Dravidian \ml{തട}, \textquotedblleft{}resistance, what impedes or stops\textquotedblright{} (DEDR 3031); the ending -\ml{സ്സം} is not settled.
\end{cousinweb}

\begin{grammarlens}[title={What you've built}]
One more word for saying what happens if.
\end{grammarlens}

\subsection*{Your turn}
\begin{itemize}
\raggedright
  \item \emph{Say it:} \emph{taṭassaṁ}
  \item \emph{Say it:} \emph{taṭassaṁ}, once more
  \item \emph{Say it:} \emph{maḻa kaḷikkŭ taṭassamāyi.} — The rain became an obstacle to the game.
  \item \emph{Recall:} say the Malayalam for a crisis, then the Malayalam for a challenge, then say \emph{taṭassaṁ} again
\end{itemize}

\begin{tcolorbox}[breakable,colback=teal!4,colframe=teal!35!black,title={Before you move on}]
What does \ml{തടസ്സം} mean? (\textquotedblleft{}An obstacle\textquotedblright{}.) Say it once more.
\end{tcolorbox}
\section[\texorpdfstring{\ml{വിശകലനം} \ml{ചെയ്യുക} (viśakalanaṁ ceyyuka)}{വിശകലനം ചെയ്യുക (viśakalanaṁ ceyyuka)}]{\ml{വിശകലനം} \ml{ചെയ്യുക} (viśakalanaṁ ceyyuka) — to analyse}

\label{lesson:ML-C448-visakalanam-ceyyuka}



\begin{quote}
Before the new one: say the Malayalam for a challenge, then the Malayalam for an obstacle.
\end{quote}

\subsection*{You'll want to know: \ml{വിശകലനം} \ml{ചെയ്യുക}}
\textbf{\ml{വിശകലനം} \ml{ചെയ്യുക}} — \emph{viśakalanaṁ ceyyuka} — \textquotedblleft{}to analyse\textquotedblright{}.

\begin{quote}
\ml{അവർ} \ml{പ്രതിസന്ധിയുടെ} \ml{കാരണങ്ങൾ} \ml{വിശകലനം} \ml{ചെയ്തു}.
\end{quote}

\begin{quote}
\emph{avar pratisandhiyuṭe kāraṇaṅṅaḷ viśakalanaṁ ceytu.}
\end{quote}

\begin{quote}
They analysed the causes of the crisis.
\end{quote}

Past: \ml{വിശകലനം} \ml{ചെയ്തു}, \emph{viśakalanaṁ ceytu}.

The noun \ml{വിശകലനം} is an analysis.

\begin{cousinweb}
\textbf{\ml{വിശകലനം} \ml{ചെയ്യുക}} (\emph{viśakalanaṁ ceyyuka}) is a phrase verb built on \ml{വിശകലനം}, \textquotedblleft{}an analysis\textquotedblright{}, a learned noun formed on the Sanskrit participle \emph{viśakalita}, \textquotedblleft{}separated, divided; sifted\textquotedblright{} (Monier-Williams), with \ml{ചെയ്യുക}, \textquotedblleft{}to do\textquotedblright{}: to take a thing apart.
\end{cousinweb}

\begin{grammarlens}[title={What you've built}]
One more word for saying what happens if.
\end{grammarlens}

\subsection*{Your turn}
\begin{itemize}
\raggedright
  \item \emph{Say it:} \emph{viśakalanaṁ ceyyuka}
  \item \emph{Say it:} \emph{viśakalanaṁ ceyyuka}, once more
  \item \emph{Say it:} \emph{avar pratisandhiyuṭe kāraṇaṅṅaḷ viśakalanaṁ ceytu.} — They analysed the causes of the crisis.
  \item \emph{Recall:} say the Malayalam for a challenge, then the Malayalam for an obstacle, then say \emph{viśakalanaṁ ceyyuka} again
\end{itemize}

\begin{tcolorbox}[breakable,colback=teal!4,colframe=teal!35!black,title={Before you move on}]
What does \ml{വിശകലനം} \ml{ചെയ്യുക} mean? (\textquotedblleft{}To analyse\textquotedblright{}.) Say it once more.
\end{tcolorbox}
\section[\texorpdfstring{\ml{രീതി} (rīti)}{രീതി (rīti)}]{\ml{രീതി} (rīti) — a method; a manner}

\label{lesson:ML-C448-riti}



\begin{quote}
Before the new one: say the Malayalam for an obstacle, then the Malayalam for to analyse.
\end{quote}

\subsection*{You'll want to know: \ml{രീതി}}
\textbf{\ml{രീതി}} — \emph{rīti} — \textquotedblleft{}a method; a manner\textquotedblright{}.

\begin{quote}
\ml{ഓരോരുത്തർക്കും} \ml{പഠിക്കാൻ} \ml{ഓരോ} \ml{രീതിയുണ്ട്}.
\end{quote}

\begin{quote}
\emph{ōrōruttarkkuṁ paṭhikkān ōrō rītiyuṇṭŭ.}
\end{quote}

\begin{quote}
Everyone has their own way of learning.
\end{quote}

In that way: \ml{ആ} \ml{രീതിയിൽ}.

\begin{cousinweb}
\textbf{\ml{രീതി}} (\emph{rīti}) is borrowed from Sanskrit \emph{rīti}, \textquotedblleft{}a going, a course; a stream\textquotedblright{} (CDIAL 10751; Monier-Williams). A method is the path a thing follows.
\end{cousinweb}

\begin{grammarlens}[title={What you've built}]
One more word for saying what happens if.
\end{grammarlens}

\subsection*{Your turn}
\begin{itemize}
\raggedright
  \item \emph{Say it:} \emph{rīti}
  \item \emph{Say it:} \emph{rīti}, once more
  \item \emph{Say it:} \emph{ōrōruttarkkuṁ paṭhikkān ōrō rītiyuṇṭŭ.} — Everyone has their own way of learning.
  \item \emph{Recall:} say the Malayalam for an obstacle, then the Malayalam for to analyse, then say \emph{rīti} again
\end{itemize}

\begin{tcolorbox}[breakable,colback=teal!4,colframe=teal!35!black,title={Before you move on}]
What does \ml{രീതി} mean? (\textquotedblleft{}A method; a manner\textquotedblright{}.) Say it once more.
\end{tcolorbox}
